\section{The same separation for pure NAND circuits}

The separation also holds when every gate is a two-input NAND. Size
means the number of NAND gates; intermediate results and input wires may
be reused without restriction. Constants may be allowed or omitted.
Negations and complemented edges are not free operations in this model.

\begin{lemma}[A positive certificate for a NAND circuit]
\label{lem:occam:nand-certificate}
Let a NAND circuit with $s\ge1$ gates compute $h$, and suppose
$h(1^d)=1$. There is a set $R\subseteq[d]$ with
$|R|\le\lfloor s/2\rfloor+1$ such that $h(x)=1$ whenever
$x_i=1$ for every $i\in R$. In particular, $h(1_R)=1$.
\end{lemma}

\begin{proof}
Evaluate every gate at $1^d$, calling it true or false according to its
value. Starting at the true output gate, retain one false-valued input
of each selected true NAND gate, and both true-valued inputs of each
selected false NAND gate. Recurse to the selected sources, stopping at
variables and constants. Identify repeated occurrences of the same
source or gate, but keep repeated edges when both input slots select
the same source. The resulting proof graph is a connected rooted
directed acyclic multigraph.

Let $T,F,L$ be its numbers of true gates, false gates, and distinct
source leaves, respectively. Every selected false gate has a selected
true parent, whereas every true gate selects only one input. Thus
$F\le T$: choose one true parent of each false gate; distinct false
gates cannot choose the same parent. There are $T+2F$ selected edges
and $T+F+L$ vertices. Connectedness gives
\[
 T+F+L-1\le T+2F,\qquad L\le F+1.
\]
Since $T+F\le s$ and $F\le T$, we obtain
$L\le\lfloor s/2\rfloor+1$.

Let $R$ consist of the distinct variable leaves, so $|R|\le L$.
Every one of these variables is one at $1^d$. Fixing them to one
forces every selected gate to retain its value at $1^d$, by induction
in circuit order: a selected false NAND has both inputs forced true,
and a selected true NAND has one input forced false. Constants retain
their values without any restriction on the input. The output
therefore remains one. The argument explicitly permits sharing and
repeated input wires.
\end{proof}

\begin{corollary}[Exact NAND complexity of conjunctions]
\label{cor:occam:nand-and}
For $k\ge2$, the conjunction of $k$ distinct variables has minimum
NAND-circuit size $2k-2$.
\end{corollary}

\begin{proof}
An AND of two signals is a NAND followed by a NAND with its two inputs
tied together. Replacing every gate in an AND tree gives $2k-2$ gates.
Conversely, every positive certificate for a conjunction must include
all its $k$ variables. The lemma gives
$k\le\lfloor s/2\rfloor+1$, hence $s\ge2k-2$.
\end{proof}

\begin{theorem}[Exponential mistakes for smallest NAND circuits]
\label{thm:occam:nand-ladder}
Fix $d\ge3$ and let $f(x)=\bigwedge_{i=1}^d x_i$. Consider any
deterministic learner which, before receiving the next input, chooses
a minimum-size NAND circuit consistent with all preceding examples.
There is a permutation of $\{0,1\}^d$, labeled throughout by $f$,
on which the learner makes at least
\[
 2^d-d-2
   =2^{s/2+1}-s/2-3
 \quad\hbox{mistakes},\qquad s=C_{\mathrm{NAND}}(f)=2d-2.
\]
The permutation may depend on the deterministic learner, but can be
fixed before the sequence is presented.
\end{theorem}

\begin{proof}
First present $1^d$ with label one, then every input of weight zero
or one with label zero. For $k=2,\ldots,d-1$, present the weight-$k$
inputs as follows.

At any point in stage $k$, every input of weight below $k$ has already
been labeled zero. If a consistent circuit has $r$ gates, the lemma
provides a positive certificate $R$ with
$|R|\le\lfloor r/2\rfloor+1$. Consistency requires $|R|\ge k$,
since otherwise $h(1_R)=1$ contradicts a previous label. Therefore
every consistent circuit has at least $2k-2$ gates.

As long as some weight-$k$ input $1_S$ remains unshown, the conjunction
$\bigwedge_{i\in S}x_i$ is consistent: it is one at $1^d$, zero below
weight $k$, and among weight-$k$ inputs it is one only at $1_S$.
It has $2k-2$ NAND gates. Thus the learner's minimum consistent size
is exactly $2k-2$. A positive certificate $R$ for its chosen circuit
has $|R|\le k$, and consistency forces $|R|\ge k$. Present $1_R$.
This input has not been shown before, the learner predicts one, and
the target label is zero because $|R|=k<d$.

Repeating this step presents every weight-$k$ input as a mistake.
Summing over $k=2,\ldots,d-1$ gives
$\sum_{k=2}^{d-1}\binom dk=2^d-d-2$. The initial points and these
levels comprise the whole hypercube. Because the learner is
deterministic, the constructed sequence can be simulated and fixed
in advance. The target labels never depend on the learner.
\end{proof}

\paragraph{A simple fixed tie rule and input order.}
For every $k\ge2$, fix the usual left-associated AND chain on a sorted
list of $k$ variable indices and replace each AND by its two-NAND
implementation. At size $2k-2$, place these canonical circuits first,
ordered lexicographically by their supports, and place all remaining
circuit descriptions afterward in any fixed order. At other sizes,
use any fixed order. Present $1^d$, then $0^d$, then the remaining inputs in
increasing Hamming weight, using lexicographic support order within
each level. At every stage $k\ge2$, the minimum consistent size is
$2k-2$, and the first surviving canonical conjunction has exactly the
support of the next input. This fixed tie rule and fixed input order
therefore realize the lower bound. No classification of all
minimum-size NAND hypotheses as conjunctions is required.

\paragraph{Conventions.}
The argument starts at $k=2$, so it applies both when outputs must be
gates and when direct wires or constants are allowed as zero-gate
outputs. If desired, the certificate lemma extends to those
zero-gate cases: a projection has a one-variable certificate and
constant one has the empty certificate. The result concerns gate
count. In a literal gate-list code with input addresses, the target's
description length is $\Theta(d\log d)$, so the lower bound becomes
$2^{\Omega(\ell/\log\ell)}$ in that description length.
